\chapter{Formelle Definition von Markov-Ketten}\label{ch:markov-formal}

Formal definieren wir eine Markov-Kette über einen diskreten Zustandsraum $S$ und eine Übergangsmatrix $P$.

\begin{mydefinition}[Markov-Kette]
    Eine \textbf{Markov-Kette} ist ein stochastischer Prozess $\{X_n : n \in \mathbb{N}_0\}$ mit Werten in einem abzählbaren Zustandsraum $S$, der die Markov-Eigenschaft erfüllt:
    \[
        P(X_{n+1} = j | X_n = i, X_{n-1} = i_{n-1}, \ldots, X_0 = i_0) = P(X_{n+1} = j | X_n = i)
    \]

    Die Wahrscheinlichkeit $p_{ij} = P(X_{n+1} = j | X_n = i)$ bezeichnet die Übergangswahrscheinlichkeit von Zustand $i$ zu Zustand $j$. Die Matrix $P = (p_{ij})_{i,j \in S}$ heisst \textbf{Übergangsmatrix} der Markov-Kette.

    Eine Markov-Kette heisst \textbf{homogen}, wenn die Übergangswahrscheinlichkeiten nicht vom Zeitpunkt $n$ abhängen, sondern nur von den beteiligten Zuständen $i$ und $j$.
\end{mydefinition}

Der Zustandsvektor $\pi^{(n)} = (\pi_1^{(n)}, \pi_2^{(n)}, \ldots)$ gibt die Wahrscheinlichkeitsverteilung der Zustände zum Zeitpunkt $n$ an, wobei $\pi_i^{(n)} = P(X_n = i)$. Der Anfangszustand wird durch den Vektor $\pi^{(0)}$ beschrieben.

Die zeitliche Entwicklung einer Markov-Kette lässt sich durch wiederholte Multiplikation mit der Übergangsmatrix $P$ berechnen:
\[
    \pi^{(n+1)} = \pi^{(n)} \cdot P
\]

Das bedeutet, dass die Verteilung nach $n$ Schritten durch $\pi^{(n)} = \pi^{(0)} \cdot P^n$ gegeben ist.